\subsection{Fibonacci en Tiempo Logarítmico (Fast Doubling)}
\label{alg:7-8}

\graybold{Preguntas guía:}

\begin{itemize}
\item ¿Me piden un término de Fibonacci (o de otra recurrencia lineal) con un índice enorme, \ten{18} o más?
\item ¿El resultado se pide módulo algo, porque el valor exacto tendría cientos de millones de dígitos?
\item ¿Puedo saltar de $n$ a $2n$ en vez de avanzar de a un término?
\item ¿La recurrencia es lineal con coeficientes constantes, de modo que también serviría exponenciación de matrices?
\end{itemize}

\graybold{Utilidad:} Calcula $F_n$ en \bigO{\log n} operaciones usando dos identidades que permiten duplicar el índice. Es la versión más rápida y corta para Fibonacci; la técnica general detrás (exponenciación binaria aplicada a una recurrencia) sirve para cualquier recurrencia lineal con coeficientes constantes, y también para contar caminos de longitud $n$ en un grafo o para aplicar $n$ veces una transformación lineal.

\graybold{Complejidad:} \bigO{\log n} multiplicaciones modulares: una por cada bit de $n$, es decir a lo sumo 60 para $n \le 10^{18}$.

\graybold{Fundamento teórico:} De la identidad de Fibonacci $F_{m+k} = F_m F_{k+1} + F_{m-1} F_k$ se derivan, tomando $m = k$, las dos fórmulas de duplicación:
\[F_{2k} = F_k\,(2F_{k+1} - F_k), \qquad F_{2k+1} = F_k^2 + F_{k+1}^2.\]
Ambas expresan el par $(F_{2k}, F_{2k+1})$ en función del par $(F_k, F_{k+1})$, así que manteniendo siempre un par consecutivo se puede duplicar el índice con cuatro multiplicaciones. Para llegar a un $n$ arbitrario se recorren sus bits de mayor a menor empezando en $(F_0, F_1) = (0,1)$: en cada paso se duplica el índice acumulado y, si el bit vale 1, se avanza además un término, pasando de $(F_{2k}, F_{2k+1})$ a $(F_{2k+1}, F_{2k+2})$. Esto es exponenciación binaria disfrazada: el índice se construye bit a bit igual que el exponente en una potenciación rápida. Como todo se hace módulo $\ten{9}+7$, las magnitudes se mantienen chicas; la única sutileza es que $2F_{k+1} - F_k$ puede quedar negativo tras reducir, y el operador módulo de Python ya devuelve un representante no negativo, así que no hace falta corregirlo a mano (en C++ sí haría falta). La alternativa clásica es elevar a la potencia $n$ la matriz de $2\times2$ cuya primera fila es $(1,\,1)$ y cuya segunda fila es $(1,\,0)$: da el mismo resultado con ocho multiplicaciones por paso en lugar de cuatro, pero tiene la ventaja de generalizarse a cualquier recurrencia lineal cambiando la matriz.~\cite{knuth1997}

\graybold{Ejercicio aplicado}

(CSES 1722 — Fibonacci Numbers) Dado $n$ con $0 \le n \le 10^{18}$, imprimir $F_n$ módulo $\ten{9}+7$.

\graybold{I/O:} Una sola línea con un entero que llega a \ten{18}, sin problema para los enteros de Python; se lee todo de una y se convierte directamente. Verificado en 500 valores (todos los de 0 a 299 y 200 al azar hasta $2\cdot10^5$) contra la iteración directa de la recurrencia, y en 65 valores gigantes (incluidos \ten{18}, $\ten{18}-1$ y \pow{2}{60}) contra una implementación independiente por exponenciación de matrices, sin discrepancias. El tiempo medido con $n = 10^{18}$ es \aprox{}0.009s, que es esencialmente el arranque del intérprete.

\graybold{Código solución}

\begin{lstlisting}[language=Python]
import sys

def solve():
    n = int(sys.stdin.buffer.read())
    MOD = 10**9 + 7
    # fast doubling: se recorren los bits de n de mayor a menor,
    # manteniendo siempre el par (F(k), F(k+1))
    a, b = 0, 1                                # (F(0), F(1))
    for bit in bin(n)[2:]:
        # de (F(k), F(k+1)) a (F(2k), F(2k+1))
        c = a * ((2*b - a) % MOD) % MOD        # F(2k)   = F(k) * (2*F(k+1) - F(k))
        d = (a*a + b*b) % MOD                  # F(2k+1) = F(k)^2 + F(k+1)^2
        if bit == "1":
            a, b = d, (c + d) % MOD            # bit 1: ademas se avanza un termino
        else:
            a, b = c, d
    print(a % MOD)

solve()
\end{lstlisting}

\graybold{Extras: exponenciación de matrices} — más lenta que fast doubling para Fibonacci, pero sirve para CUALQUIER recurrencia lineal con coeficientes constantes: basta cambiar la matriz. Verificada contra la solución principal en 65 valores hasta \ten{18}.

\begin{lstlisting}[language=Python]
MOD = 10**9 + 7

def matpow_fib(n):
    def mul(X, Y):
        return [[(X[0][0]*Y[0][0] + X[0][1]*Y[1][0]) % MOD,
                 (X[0][0]*Y[0][1] + X[0][1]*Y[1][1]) % MOD],
                [(X[1][0]*Y[0][0] + X[1][1]*Y[1][0]) % MOD,
                 (X[1][0]*Y[0][1] + X[1][1]*Y[1][1]) % MOD]]
    R = [[1,0],[0,1]]          # identidad
    B = [[1,1],[1,0]]          # matriz de la recurrencia de Fibonacci
    while n:                   # exponenciacion binaria
        if n & 1: R = mul(R, B)
        B = mul(B, B); n >>= 1
    return R[0][1]             # R = B^n, y B^n[0][1] = F(n)

# para otra recurrencia lineal a_n = c1*a_{n-1} + ... + ck*a_{n-k}, la matriz es
# la fila [c1, c2, ..., ck] arriba y la identidad desplazada debajo; el vector
# inicial son los k primeros terminos.
\end{lstlisting}
